\section[Demonstrations]{\secstyle{Demonstrations:}}

\subsection[Ensembles \& majorations]{\subsecstyle{\textbf{Ensembles \& majorations:}}}
Soit A et B des parties non vides et bornées de \(\R\) telles que \(A \cap B\) est non vide.
\subsubsection*{\subsubsecstyle{(1) Montrer que si \(A \subseteq B \text{ alors } \sup(A) \leq \sup(B)\)}}
\begin{proof}[\unskip\nopunct]
   Supposons \(A \subseteq B\), soit \(a \in A\) et \(b \in B\) on sait que:
   \[
      b \leq \sup(B)    
   \]
   Alors en particulier on a:
   \[
      a \leq \sup(B) \;\; \text{(Car \(A \subseteq B\))} 
   \]
   Par passage à la borne supérieure:
   \[
      \sup(A) \leq \sup(B)    
   \]
\end{proof}

\subsubsection*{\subsubsecstyle{(2) (a) Montrer que \(A \cup B\) admet une borne supérieure finie}}
Soit \(\alpha \in A \cup B\), alors par définition on a 3 cas à traiter:
\[
   \begin{cases}
      \alpha \in A \implies \alpha \leq \sup(A)\\
      \alpha \in B \implies \alpha \leq \sup(B)\\
      \alpha \in A \cap B \implies \alpha \leq \sup(B) \text{ et }\alpha \leq \sup(A) \text{ (Car \(\alpha \in A\) et \(\alpha \in B\))}\\
   \end{cases}
\]
On remarque que dans chaque cas, \(\alpha\) est majoré par une quantité finie, et ceci étant vrai pour un \(\alpha\) quelconque de \(A \cup B\), alors \(A \cup B\) est majorée et donc on conclut par la propriété de \(\R\).

\subsubsection*{\subsubsecstyle{(2) (b) Montrer que \(\sup(A \cup B)= \max(\sup(A), \sup(B))\)}}
Posons \(M := \max(\sup(A), \sup(B))\), prouvons l'égalité par double inégalité:

\begin{proof}[\unskip\nopunct]
   \boxed{\leq} Soit \(\alpha\) un élément quelconque de \(A \cup B\), on remarque tout d'abord que:
   \[
      (\sup(A) \leq M \text{ et } \sup(B) \leq M) \implies \alpha \leq M
   \]
   Par passage la borne supérieure on en conclut que \(\sup(A \cup B)\leq M \).\<

   \boxed{\geq} Montrons que:
   \[
      \begin{cases}
         \sup(A \cup B) \geq \sup(A)\\
         \sup(A \cup B) \geq \sup(B)
      \end{cases} \text{ et donc que: \(\sup(A \cup B) \geq M\)}
   \]
   Soit \(\alpha \in A\) (resp. \(\alpha \in B\)), alors \(\alpha \in A \cup B\), et donc en particulier:
   \[
      \sup(A \cup B) \geq \alpha 
   \]
   Par passage à la borne supérieure on trouve respectivement:

   \begin{alignat}{2}
      \sup(A \cup B) &\geq \sup(A) \\  
      \sup(A \cup B) &\geq \sup(B)   
   \end{alignat}
   Donc \(\sup(A \cup B) \geq M\) 

\end{proof}

\subsubsection*{\subsubsecstyle{(2) (c) Montrer que \(\inf(A \cup B)= \min(\inf(A), \inf(B))\)}}
Posons \(m := \min(\inf(A), \inf(B))\), prouvons l'égalité par double inégalité:

\begin{proof}[\unskip\nopunct]
   \boxed{\geq} Soit \(\alpha \in A \cup B\), on remarque tout d'abord que:
   \[
      (\inf(A) \geq m \text{ et } \inf(B) \geq m ) \implies \alpha \geq m
   \]
   Par passage à la borne inférieure, on conclut que:
   \[
      \inf(A \cup B) \geq m    
   \]
   \boxed{\leq} Montrons que:
   \[
      \begin{cases}
         \inf(A \cup B) \leq \inf(A) \\
         \inf(A \cup B) \leq \inf(B)
      \end{cases} \text{ et donc que: \(\inf(A \cup B) \leq m\)}
   \]
   Soit \(\alpha \in A\) (resp. \(\alpha \in B\)), alors \(\alpha \in A \cup B\) et donc:
   \[
         \inf(A \cup B) \leq \alpha
   \]
   Par passage à la borne inférieure, on trouve respectivement:
   \begin{alignat*}{2}
      \inf(A \cup B) \leq \inf(A)\\
      \inf(A \cup B) \leq \inf(B)
   \end{alignat*}
   Donc \( \inf(A \cup B) \leq m\)
\end{proof}

\subsubsection*{\subsubsecstyle{(3) (a) Montrer que \(\sup(A \cap B)\) admet une borne supérieure finie}}

Soit \(\alpha \in A \cap B\), A et B sont bornées, donc en particulier A et B sont majorées, en particulier on a:
\[
   \exists M, \forall a \in A \; ; \; M \geq a
\]   

Soit \(M\) le majorant de \(A\), on sait que \(\alpha \in A\) donc \(M \geq \alpha\) et donc \(A \cap B\) est majorée, on conclut par la propriété de \(\R\).

\subsubsection*{\subsubsecstyle{(3) (b) Montrer que \(\sup(A \cap B) \leq \min(\sup(A), \sup(B))\)}}
\begin{proof}[\unskip\nopunct]
   Soit \(\alpha \in A \cap B\), montrons que \(\alpha \leq \sup(A)\) et \(\alpha \leq \sup(B)\) et par suite que \(\alpha \leq \min(\sup(A), \sup(B))\). \<

   \(\alpha \in A \cap B \Longleftrightarrow \alpha \in A \text{ et } \alpha \in B\), donc on en déduis que:
   \[
      \alpha \leq \sup(A) \text{ et } \alpha \leq \sup(B) \implies \alpha \leq \min(\sup(A), \sup(B))\\
   \]

   Par passage à la borne supérieure on en déduis que \(\sup(A \cap B) \leq \min(\sup(A), \sup(B))\)
\end{proof}

\underline{Exemple tel que \(\sup(A \cap B) \neq \min(\sup(A), \sup(B)\)):}\<

Soit \(I\) l'intervalle \(\icc{0}{5}\) et \(E\):= \(\bigl\{ 0, 5\bigl\}\)  et donc \(I \cap E = \bigl\{ 0 \bigl\}\).

Alors on a \(\sup(I) = 5\) et \(\sup(E) = 5\), et donc \(\min(\sup(I), \sup(E)) = 5\). 

Mais on voit que \(\sup(I \cap E) = 0 \neq 5\).

\subsubsection*{\subsubsecstyle{(3) (c) Montrer que \(\inf(A \cap B) \geq \max(\inf(A), \inf(B))\)}}
\begin{proof}[\unskip\nopunct]
   Soit \(\alpha \in A \cap B\), montrons que \(\alpha \geq \inf(A)\) et \(\alpha \geq \inf(B)\), et par suite que \(\alpha \geq \max(\inf(A), \inf(B))\).\<

   \(\alpha \in A\) et \(\alpha \in B\) donc on a:
   \[
      \alpha \geq \inf(A) \text{ et } \alpha \geq \inf(A) \implies \alpha \geq \max(\inf(A), \inf(B)) 
   \]
   On conclut par passage à la borne inférieure.
\end{proof} 

\underline{Exemple tel que \(\inf(A \cap B) \neq \max(\inf(A), \inf(B)\)):}\<

Soit \(I\) l'intervalle \(\icc{3}{5}\) et \(E\):= \(\bigl\{ 0, 5\bigl\}\)  et donc \(I \cap E = \bigl\{ 5 \bigl\}\).

Alors on a \(\inf(I) = 3\) et \(\inf(E) = 0\), et donc \(\max(\inf(I), \inf(E)) = 3\). 

Mais on voit que \(\inf(I \cap E) = 5 \neq 3\).
\pagebreak

\subsection[Ecart d'un ensemble]{\subsecstyle{\textbf{Ecart d'un ensemble:}}}
Soient \(A \subseteq \R \) une partie bornée, et B l'ensemble:
\[
B:= \Bigl\{ |x - y| \; ; \; x, y \in A \Bigl\}
\]
\subsubsection*{\subsubsecstyle{(1) Montrer que B admet une borne supérieure finie \(\delta(A)\)}}
\begin{proof}[\unskip\nopunct]
On sait que A est bornée, ie:
\begin{align*}
   \exists M \in \R\; , \; &\forall x \in A \; ; \; M \geq x \\
   \exists m \in \R\; , \; &\forall x \in A \; ; \; m \leq x \\
\end{align*}
Soit M et m les bornes de A. Distinguons plusieurs cas:
\[
   \begin{cases}
      \text{Si } x \geq y \text{, alors } |x - y| = x - y\\
      \text{Sinon, } |x - y| = y - x\\
   \end{cases}
\]
On cherche donc à majorer \(x - y\) et \(y - x\), on a:
\begin{align*}
   \sup(A) &\geq x\\
   -\inf(A) &\geq -y
\end{align*}
Puis par somme:
\[
      \sup(A) - \inf(A) \geq x - y 
\]

On a donc prouvé que B est majorée, et donc possède une borne supérieure (par definition de \(\R\))
Par symétrie, on trouve le même majorant pour \(y - x\). Et par une procédure analogue, on peut prouver l'existence de la borne inférieure de B. 
\end{proof}
\subsubsection*{\subsubsecstyle{(2) Prouver que \(\delta(A) = \sup(A) - \inf(A)\)}}
\begin{proof}[\unskip\nopunct]
Posons M = \(\sup(A) - \inf(A)\).
On a montré en (1) que M majorait B. Il reste à montrer que M est le plus petit des majorants, ie:
\[
   \forall \epsilon > 0, \exists x, y \in A \; ; \; M - \epsilon < x - y \text{ (La procédure est symétrique dans le second cas \(y - x\))}
\]
Soit \(\epsilon\) strictement positif, si on pose \(x \in \ioo{\sup(A) - \frac{\epsilon}{2}}{\sup(A)} \cap A \text{ et } y \in \ioo{\inf(A)}{\inf(A) + \frac{\epsilon}{2}} \cap B \text{ et donc } -y \in \ioo{-\inf(A) - \frac{\epsilon}{2}}{-\inf(A)}\cap B\), on a:
\[
   \begin{cases}
      \sup(A) - \frac{\epsilon}{2} < x < \sup(A)   \\
      -\inf{A} - \frac{\epsilon}{2} < -y < -\inf(A) \\
   \end{cases} \implies
   x - y> \sup(A) - \frac{\epsilon}{2} \inf{A} - \frac{\epsilon}{2} = M - \epsilon\\ 
\]
Aussi on a \(x \in A\) et \(y \in B\) car \(x\) est dans un intervalle qui a une intersection non-vide avec A (car inférieur au sup, et respectivement pour \(y\) qui est inférieur à l'inf, et car on sait que A et B sont non vides).

Donc M est un majorant, et tout réel plus petit que M n'est plus un majorant, et donc \(\delta(A) = \sup(A) - \inf(A)\).
\end{proof}
\pagebreak

\subsection[Ensemble somme]{\subsecstyle{\textbf{Ensemble somme:}}}
Soient A, et B des parties de \(\R\) et A + B l'ensemble:
\[
A + B:= \Bigl\{ a + b \; ; \; a\in A, b\in  B\Bigl\}
\]
\subsubsection*{\subsubsecstyle{(1) Montrer que \(\sup(A + B)\) existe:}}
\begin{proof}[\unskip\nopunct]
A et B sont bornées donc il existe une borne supérieure et inférieure de A et B (par la propriété de \(\R\)), donc on a:
\begin{align*}
   \forall x \in A \; &; \; x \leq \sup(A) \\
   \forall y \in B \; &; \; y \leq \sup(B)
\end{align*}
Donc on a \(x + y \leq \sup(A) + \sup(B)\). Or \(\sup(A) + \sup(B)\) est une quantité finie, donc on a trouvé un majorant de \(A + B\) qui est fini. Donc \(A + B\) admet une borne supérieure finie. La procédure est similaire pour \(\inf(A + B)\).
\end{proof}
\subsubsection*{\subsubsecstyle{(2) Montrer que \(\sup(A + B) = \sup(A) + \sup(B)\)}}
\begin{proof}[\unskip\nopunct]
D'après le résultat de la question (1) on sait que \(\sup(A) + \sup(B)\) majore \(A + B\), notons M ce majorant et démontrons que tout nombre plus petit ne majore plus \(A + B\), ie:
\[
   \forall \epsilon > 0, \exists (a, b) \in A \times B \; ; \; a + b > M - \epsilon   
\]
Soit \(\epsilon\) strictement positif, on cherche un couple \((a, b) \in A \times B\), tel que \(a + b > M - \epsilon\).

Posons:
\begin{alignat*}{2}
   a \in \ioo{\sup(A) - \frac{\epsilon}{4}}{\sup(A)} \cap A \\
   b \in \ioo{\sup(B) - \frac{\epsilon}{4}}{\sup(B)} \cap B
\end{alignat*}

Alors on a bien:
\[
   a + b > M - \frac{\epsilon }{2} > M - \epsilon   
\]
Aussi on sait que \(a \text{ et } b\) sont dans leurs ensembles respectifs car:
\[
   \begin{cases}
      \text{A et B sont non-vides}\\ 
      \text{Ils appartient à un intervalle majoré par leurs borne supérieures respectives}\\
      \text{Cet intervalle est intersecté avec leurs ensembles respectifs}\\    
   \end{cases}   
\]

Donc on a montré que tout nombre plus petit que \(\sup(A) + \sup(B)\) n'est plus un majorant, et donc:
\[
   \sup(A) + \sup(B) = \sup(A + B)
\]
\end{proof}

\subsubsection*{\subsubsecstyle{(3) Montrer que si A et B sont des intervalles, A + B est un intervalle:}}
Supposons A et B des intervalles et montrons que A + B est un intervalle.
\begin{proof}[\unskip\nopunct]
A et B étant des intervalles, pour \(a, b, c, d\) des réels quelconques, on a par définition:
\begin{align*}
   A &:= \Bigl\{ x \; ; \; a < x < b \Bigr\} \\
   B &:= \Bigl\{ y \; ; \; c < y < d \Bigr\}  
\end{align*}
A + B étant défini comme la somme des éléments de A avec les éléments de B, on a:
\begin{align*}
   x > a \text{ et } y > c &\implies x + y > a + c \\
   x < b \text{ et } y < d &\implies x + y < b + d
\end{align*}
Et donc:
\[
   A + B := \Bigl\{ \phi \; ; \; a + c < \phi < b + d \Bigr\}
\]
Ce qui est, par définition des intervalles, équivalent à:
\[
   A + B = \ioo{a + c}{b + d}
\]
\end{proof}
\pagebreak

\subsection[Ensemble produit]{\subsecstyle{\textbf{Ensemble produit:}}}  
Soient A, et B des parties de \(\R\) et AB l'ensemble:
\[
AB:= \Bigl\{ ab \; ; \; a\in A, b\in  B\Bigl\}
\]
\subsubsection*{\subsubsecstyle{(1) Montrer que \(AB\) admet des bornes:}}
\begin{proof}[\unskip\nopunct]
A et B sont bornées donc il existe une borne supérieure et inférieure de A et B (par la propriété de \(\R\)), donc on a:
\begin{align*}
   \forall x \in A \; &; \; \inf(A) \leq x \leq \sup(A) \\
   \forall y \in B \; &; \; \inf(B) \leq y \leq \sup(B)
\end{align*}
Posons \(\lambda_X := \max(|\sup(X)|, \;|\inf(X)|)\), alors on a cette propriété qui découle aisément des définitions:
\begin{align*}
   \forall x \in A \; ; \; -\lambda_A \leq - |\inf(A)| \leq & \;\boxed{\inf(A) \leq x \leq \sup(A)} \leq |\sup(A)| \leq \lambda_A \\
   \forall y \in B \; ; \; -\lambda_B \leq - |\inf(B)| \leq & \;\boxed{\inf(B) \leq y \leq \sup(B)} \leq |\sup(B)| \leq \lambda_B
\end{align*}
Par transitivité:
\begin{align*}
   (\forall x \in A \; &; \; -\lambda_A \leq x \leq \lambda_A) \\
   (\forall y \in B \; &; \; -\lambda_B \leq y \leq \lambda_B)
\end{align*}
Par définition de la valeur absolue, on a donc:
\begin{align*}
   \forall x \in A \; &; \; |x| \leq \lambda_A \\
   \forall y \in B \; &; \; |y| \leq \lambda_B
\end{align*}
Ensuite, car les quantités sont positives, on a (pour tout couple \(xy\) de \(A \times B\)):
\[
   |xy| \leq \lambda_A\lambda_B
\]
A et B étant bornées, \(\lambda_A\lambda_B\) est une quantité finie et donc par définition de la valeur absolue on a:
\[
   xy \in \icc{-(\lambda_A\lambda_B)}{\lambda_A\lambda_B}
\]
Ce qui signifie que \(\lambda_A\lambda_B\) majore \(xy\) et \(-(\lambda_A\lambda_B)\) minore \(xy\).\+
On viens donc de prouver \(\sup(AB)\) et \(\inf(AB)\) existent tous deux et sont finis (par définition de \(\R\)).
\end{proof}
\subsubsection*{\subsubsecstyle{(2) Montrer que \(A, B \subseteq \R^+ \implies \sup(AB) = \sup(A)\sup(B)\):}}

\begin{proof}[\unskip\nopunct]
Soit \(x \in A\), \(y \in B\) et \(\phi \in AB\), alors on a:
\begin{align*}
   &0 \leq \inf(A) \leq x \leq \sup(A) \\
   &0 \leq \inf(B) \leq y \leq \sup(B)
\end{align*}
D'aprés la question (1), \(AB\) est bornée, donc on a:
\[
   \inf(AB) \leq \phi \leq \sup(AB)
\]
Or si \(x, y\) sont positifs, alors AB ne contient que des réels positifs (règle des signes). Donc \(AB \subseteq \R^+\).

Par les opérations sur les inégalités, on a:
\[
   0 \leq xy \leq \sup(A)\sup(B)\\
\]
Ceci est vrai pour tout \(xy\) et donc pour tout \(\phi \in AB\), ie:
\[
   0 \leq \phi \leq \sup(A)\sup(B)\\
\]
Par passage à la borne supérieure, on voit donc que \(\sup(A)\sup(B)\) majore \(AB\). Notons M ce majorant.\<

Montrons que M est le plus petit des majorants, ie:
\[
   \forall \epsilon > 0, \exists (a, b) \in A \times B \; ; \; ab > M - \epsilon
\]

Posons:
\begin{alignat*}{2}
   a &\in \ioo{\sup(A) - \frac{\epsilon}{2}}{\sup(A)} \cap A \Longleftrightarrow \sup(A) - \frac{\epsilon}{2} < a < \sup(A) \text{ ET } a \in A\\
   b &\in \ioo{\sup(B) - \frac{\epsilon}{2}}{\sup(B)} \cap B \Longleftrightarrow \sup(B) - \frac{\epsilon}{2} < b < \sup(B) \text{ ET } b \in B
\end{alignat*}

On en déduis que:

\end{proof}
\subsubsection*{\subsubsecstyle{(3) Qu'en est-il si \(A\) et \(B\) contient des réels négatifs:}}
La relation de (2) n'est plus vérifiée car des contre exemples existent, supposons \(A = \icc{-5}{-2}\) et \(B = \icc{-3}{-1}\) des parties de \(\R\) qui contiennent des réels négatifs, alors on a:

\begin{alignat*}{2}
&\sup(A) = -2 \\
&\sup(B) = -1 \\
&\sup(A)\sup(B) = 2 \\
\end{alignat*}
Mais il est aisé de voir que l'ensemble \(AB\) contient par exemple l'élement \(-4 \in A\) et l'élement \(-3 \in B\) et le produit des deux (qui est par définition dans \(AB\)) est égal à \(12 > \sup(A)\sup(B)\).\<

Par ailleurs, cela découle de la preuve du (1) car on a montré que AB est majoré par \(\max(|\sup(A)|, \;|\inf(A)|) \times \max(|\sup(B)|, \;|\inf(B)|)\).

\subsubsection*{\subsubsecstyle{(4) Montrer que si A et B sont des intervalles alors AB est un intervalle:}}

(a faire).
\pagebreak

\section[Exercices]{\secstyle{Exercices:}}

\subsection[Etude d'ensembles]{\subsecstyle{\textbf{Etude d'ensembles:}}}
\[
   \E:= \Bigl\{ 1 + n^3 \; ; \; n \in \N \Bigl\}
\]

\subsubsection*{\subsubsecstyle{Cherchons un intervalle qui contient précisément \(\E\):}}

Soit \(\alpha\) un élément de \(\E\).\+
On remarque que \(n^3\) étant croissante, avec \(0 \leq n^3\), on a \(1 \leq \alpha\), on en conclus que \(\E \subseteq \ico{1}{+\infty}\)

\subsubsection*{\subsubsecstyle{Cherchons l'ensemble des minorants de \(\E\):}}
On sait que \(1 \leq \alpha\), donc on conjecture que l'ensemble des minorants est l'ensemble \(E := \ioc{-\infty}{1}\), démontrons que \(MINO(\E) = E\):
\begin{proof}[\unskip\nopunct]
   Soit \(m \in E\), prouvons que \(m \in MINO(\E)\). On sait que \(1 \leq \alpha\) et que \(m \leq 1\) donc par transitivité, on peut conclure que \(m \in MINO(\E)\), car \(\alpha\) est un élément quelconque de \(\E\).\<
   
   Montrons maintenant que tout nombre qui n'est pas dans E n'est pas un minorant, ie:
   \[
      \forall x > 1, \exists \lambda \in \E \; ; \; \lambda < x
   \]
   Soit \(x > 1\) un réel, posons \(\lambda\) l'élément de \(\E\) pour \(n = 0\) alors, on a:
   \[
      \lambda = 1 < x 
   \]
   Donc tout nombre qui n'est pas dans E n'est pas un minorant, on a donc prouvé que \(MINO(\E) = E\)
\end{proof}

\subsubsection*{\subsubsecstyle{Cherchons l'ensemble des majorants de \(\E\):}}
Démontrons que \(\E\) n'est pas majoré, ie:
\[
   \forall m \in \R, \exists \lambda \in \E \; ; \; \lambda > m   
\]
\begin{proof}[\unskip\nopunct]
   Soit \(m\) un réel, on cherche un \(n \in \N\) tel que l'élément de \(\E\) pour ce \(n\) soit supérieur à \(m\). Posons \(n:= \lfloor m \rfloor + 1\), qui est bien un entier naturel, et par suite:
   \[
      1 + n^3 > m
   \]
   Donc un tel \(n\) convient ainsi que l'élément de \(\E\) correspondant, on en conclut que l'ensemble n'est pas majoré. 
\end{proof}

\subsubsection*{\subsubsecstyle{Maximum et minimum de $\E$:}}
\(\E\) admet un minimum \(\min(\E) := 1\) car c'est un minorant qui appartient à \(\E\).\+
\(\E\) n'admet pas de maximum car \(\E\) n'est pas majoré.

\subsubsection*{\subsubsecstyle{Sup et Inf de $\E$:}}
\(\E\) admet un minimum donc \(\sup(\E) = \min(\E) = 1\).\+
\(\E\) n'admet pas de borne supérieure car \(\E\) n'est pas majoré.
\pagebreak




\[ 
   \E:= \Bigl\{ {(-1)}^n + \frac{1}{n^2} \; ; \; n\in\N^* \Bigl\}
\]

\subsubsection*{\subsubsecstyle{Cherchons un intervalle qui contient précisément \(\E\):}}
Soit \(\alpha\) un élément de \(\E\). Etudions les différents cas selon la parité de $n$:

\[
   \begin{cases}
      \text{Si n est pair on a:}&      
         \begin{cases}
            {(-1)}^n = 1 \\
            0 < \frac{1}{n^2} \leq \frac{1}{4}
         \end{cases} \implies 1 < \alpha \leq \frac{5}{4} \\\\
      \text{Si n est impair:}&             
         \begin{cases}
            {(-1)}^n = -1 \\
            0 < \frac{1}{n^2} \leq 1
         \end{cases} \implies -1 < \alpha \leq 0\\
   \end{cases}
\]
Donc dans le cas général, on trouve que \(-1 < \alpha \leq \frac{5}{4}\).

\subsubsection*{\subsubsecstyle{Cherchons l'ensemble des minorants de $\E$:}}

On conjecture que $\ioc{-\infty}{-1}$ = $MINO(\E)$, prouvons le par double inclusion:

\begin{proof}[\unskip\nopunct]
   \underline{Sens direct:}\+
   Soit $m \in \ioc{-\infty}{-1}$, par définition $m \leq -1$, or on sait que $-1 < \alpha$ et donc on a $m \leq -1 < \alpha$ et donc \(m\) est bien un minorant de $\E$. (Ceci étant vrai pour un \(\alpha \in \E\) quelconque, on a bien validé la définition de \(m\) minorant de $\E$).\<
   
   \underline{Sens indirect:}\+
   Soit $m$ un élément de $MINO(\E)$, prouvons que \(m\) est nécessairement dans $\ioc{-\infty}{-1}$ par contraposition, ie prouvons:
   \[
      m \notin \ioc{-\infty}{-1} \implies m \notin MINO(\E)
   \]
   Supposons $m > -1$, montrons qu'il existe un $\alpha \in \E$ tel que $m > \alpha$ et donc que si \(m\) n'est pas dans l'intervalle, \(m\) n'est pas un minorant.\<

   \underline{Analyse:}
   On cherche donc un \(n \in \N^*\) tel que:
   \[
      {(-1)}^n + \frac{1}{n^2} < m
   \]
   Cherchons un n qui convient qui serait impair (ce qui facilitera le calcul) alors on a:

   \begin{alignat}{2}
      && \;\; \frac{1}{n^2} &< m + 1\\
      \Longleftrightarrow&& \;\; n^2 &> \frac{1}{m+1}\\
      \Longleftrightarrow&& \;\; n &> \frac{1}{\sqrt{m+1}}
   \end{alignat}
   (1 et 2): Car \(n^2 \text{ et }m + 1 \geq 0\)\<
   
   Donc on cherche un \(n\) entier tel que \(n > \frac{1}{\sqrt{m+1}}\). On en déduis que les deux entier \(\lfloor\frac{1}{\sqrt{m+1}}\rfloor + 1\) et \(\lfloor\frac{1}{\sqrt{m+1}}\rfloor + 2\) satisfont l'inégalité et l'un des deux est bien impair.\<

   \underline{Synthèse:}
   Posons $n$ l'entier impair parmi:
   \[
      \lfloor\frac{1}{\sqrt{m+1}}\rfloor + 1 \text{ et } \lfloor\frac{1}{\sqrt{m+1}}\rfloor + 2 
   \]
   Alors (par l'analyse précédente), on a:
   \begin{alignat}{2}
      && \;\; n > \frac{1}{\sqrt{m+1}} \\
      \implies&& {(-1)}^n + \frac{1}{n^2} &< m
   \end{alignat}
   On a bien trouvé un \(n\) entier, et l'élement de \(\E\) construit avec ce n montre bien que \(m\) n'est pas un minorant, et donc si un minorant existe, il est forcément dans l'intervalle $\ioc{-\infty}{-1}$.

\end{proof}

\subsubsection*{\subsubsecstyle{Cherchons l'ensemble des majorants de \(\E\):}}

On conjecture que \(\ico{\frac{5}{4}}{+\infty} = MAJO(\E) \), prouvons le par double inclusion:

\begin{proof}[\unskip\nopunct]
   \underline{Sens direct:}\+
   Soit \(M \in \ico{\frac{5}{4}}{+\infty}\), par definition \(M \geq \frac{5}{4}\), or on sait que \(\frac{5}{4} \geq \alpha\), donc \(M \geq \alpha\), en d'autres termes, \(M\) est bien un majorant de \(\E\). (Ceci étant vrai pour un \(\alpha \in \E\) quelconque, on a bien validé la définition de \(M\) majorant de $\E$)\<
   
   \underline{Sens indirect:}\+
   Soit $M$ un élément de $MAJO(\E)$, prouvons que $M$ est nécessairement dans l'intervalle $\ico{\frac{5}{4}}{+\infty}$ par contraposition:\<
   Suppposons $M \notin \ico{\frac{5}{4}}{+\infty}$, ie $M < \frac{5}{4}$, et montrons qu'il existe un $\alpha \in \E$ tel que $M < \alpha$ et donc que si \(M\) n'est pas dans l'intervalle, \(M\) n'est pas un majorant.\<

   On cherche donc un \(n\) tel que:
   \[
      {(-1)}^n + \frac{1}{n^2} > M
   \]
   Posons \(\lambda\) l'élément de \(\E\) pour \(n = 2\), par un simple calcul on a:
   \[
      \lambda = {(-1)}^2 + \frac{1}{2^2} = \frac{5}{4}
   \]
   Finalement, on a bien \(\lambda = \frac{5}{4} > M\) et donc pour que \(M\) soit un majorant, il doit nécessairement être dans $\ico{\frac{5}{4}}{+\infty}$.

\end{proof}

\subsubsection*{\subsubsecstyle{Maximum et minimum de $\E$:}}

Le maximum de $\E$ est un majorant qui appartient à $\E$, les majorants étant les réels $x \geq 2$ et $2 \in \E$, alors 2 est le maximum de $\E$.

Le minimum de $\E$ est un minorant qui appartient à $\E$, les minorants étant les réels $x < -1$ et $-1 \notin \E$, alors $\E$ n'a pas de minimum.

\subsubsection*{\subsubsecstyle{Sup et Inf de $\E$:}}

sup$(\E)$ = 2 car 2 est la maximum de $\E$ et si un maximum existe, c'est aussi la borne supérieure.

inf$(\E)$ = -1 car max$(MINO(\E)) = -1$.\pagebreak


\[
   \E:= \ioc{0}{1} \cap \Q
\]
Soit \(\alpha \in \E\), on remarque directement que \(0 < \alpha \leq 1\).

\subsubsection*{\subsubsecstyle{Cherchons l'ensemble des majorants $MAJO(\E$) de $\E$:}}
On sait que \(\alpha \leq 1\), donc on conjecture que \(MAJO(\E) = \ico{1}{+\infty}\). Démontrons le par double inclusion:
\begin{proof}[\unskip\nopunct]
   \underline{Sens direct:}\+
   Soit \(M \in \ico{1}{+\infty}\), alors \(1 \leq M\), et on sait que \(\alpha \leq 1\), donc par transitivité, on conclut que \(M \in MAJO(\E)\), car \(\alpha\) est un élément quelconque de \(\E\).\<
   
   \underline{Sens indirect:}\+
   Soit \(M \in MAJO(\E)\), supposons que M ne soit pas dans \(\ico{1}{+\infty}\), ie \(M \in \ioo{-\infty}{1}\), et montrons que M ne peut pas être un majorant et donc que:
   \[
      \exists \lambda \in \E \; ; \; \lambda > M    
   \]
   Notons que \(M \in \ioo{-\infty}{1} \Longleftrightarrow M < 1\), or on sait que \(1 \in \E\), donc \(M\) ne peut pas être un majorant, et donc les majorants sont nécessairement dans \(\ico{1}{+\infty}\).
\end{proof}

\subsubsection*{\subsubsecstyle{Cherchons l'ensemble des minorants $MINO(\E$) de $\E$:}}
On sait que \(\alpha > 0\) donc on conjecture que \(MINO(\E) = \ioc{-\infty}{0}\). Démontrons le par double inclusion:

\begin{proof}[\unskip\nopunct]
   \underline{Sens direct:}\+
   Soit \(m \in \ioc{-\infty}{0}\), alors \(m \leq 0\), or on sait que \(0 < \alpha\), donc par transitivité on peut conclure que \(m \in MINO(\E)\) car \(\alpha\) est un élément quelconque de \(\E\).\<

   \underline{Sens indirect:}\+
   Soit \(m \in MINO(\E)\), supposons que m ne soit pas dans \(\ioc{-\infty}{0}\), ie \(m \in \ioo{0}{+\infty}\) et montrons que m ne peut pas être un minorant et donc que:
   \[
      \exists \lambda \in \E \; ; \; \lambda < m   
   \]
   Notons que \(m \in \ioo{0}{+\infty} \Longleftrightarrow m > 0\). Cherchons un \(\lambda\) de \(\E\) tel que \(\lambda < m\).\<

   Considérons deux cas:
   \[
      \begin{cases}
         0 < m \leq 1 \\
         m > 1
      \end{cases}
   \]
   Si \(m > 1\), on voit que \(1 \in \E\) convient.\+
   Si \(m \in \ioc{0}{1}\), alors on considère l'intervalle \(\ioo{0}{m}\) et on pose \(\lambda\) un rationnel de cet intervalle. Un tel rationnel existe car on sait que \(\Q\) est dense dans \(\R\).

   On en conclut que peut importe la valeur de m, si m est un minorant, il ne peut pas être dans \(\ioo{0}{+\infty}\) et donc il est nécessairement dans \(\ioc{-\infty}{0}\).
\end{proof}

\subsubsection*{\subsubsecstyle{Max et Min de $\E$:}}
\(\E\) admet un maximum \(\max(\E) = 1\) car c'est un minorant qui appartient à \(\E\).\+
\(\E\) n'admet pas de minimum car \(\E \cap MINO(\E) = \emptyset\)

\subsubsection*{\subsubsecstyle{Sup et Inf de $\E$:}}
\(\E\) admet un minimum donc \(\sup(\E) = \max(\E) = 1\).\+
\(\E\) est minorée donc on a \(\inf(\E) = \max(MINO(\E)) = 0\).\+

\pagebreak


\[ 
   \boxed{\E:= \Bigl\{ x \in \Q\; ; \; x^2 \leq 2 \Bigl\}}
\]

\subsubsection*{\subsubsecstyle{Caractérisons \(\E\):}}
Soit $\alpha$ un élément de $\E$, alors on a $\alpha^2 \leq 2 \Longleftrightarrow |\alpha| \leq \sqrt{2} \Longleftrightarrow -\sqrt{2} \leq \alpha \leq \sqrt{2}$.
($\leq$ étant une disjonction, l'inégalité est valide même pour $\alpha \in \Q$, mais on observe que $\forall \alpha \in \Q \; ; \;|\alpha| \neq \sqrt{2}$ car $\sqrt{2} \notin \Q$).\<

Plus précisément, on a $-\sqrt{2} < \alpha < \sqrt{2}$, ou aussi $\E = \ioo{-\sqrt{2}}{\sqrt{2}} \cap \Q$.

\subsubsection*{\subsubsecstyle{Cherchons l'ensemble des majorants ($MAJO(\E$)) de $\E$:}}

On conjecture que $\ico{\sqrt{2}}{\infty} = MAJO(\E)$, prouvons le par double inclusion:

\begin{proof}[\unskip\nopunct]
   \underline{Sens direct:}\+
   Soit $M \in \ico{\sqrt{2}}{\infty}$, ie $M \geq \sqrt{2}$, or on sait que $\sqrt{2} > \alpha$, donc $M \geq \alpha$, en d'autres termes, $M$ est bien un majorant de $\E$. (Ceci étant vrai pour un \(\alpha \in \E\) quelconque, on a bien validé la définition de \(M\) majorant de $\E$)\<

   \underline{Sens indirect:}\+
   Soit $M$ un élément de $MAJO(\E)$, prouvons que $M$ est nécessairement dans l'intervalle $\ico{\sqrt{2}}{\infty}$ par contraposition:\\
   Soit $M \notin \ico{\sqrt{2}}{\infty}$, ie $M < \sqrt{2}$

   Montrons qu'il existe un $\alpha \in \E$ tel que $\alpha > M$, ie que $M$ n'est pas un majorant.\< 

   Considérons l'intervalle $\icc{M}{\sqrt{2}}$, on cherche un $\alpha \in \Q$ tel que $M < \alpha < \sqrt{2}$.

   Entre chaque réel il est possible de trouver un rationnel (car $\Q$ est dense dans $\R$), donc posons $\alpha$ un rationnel de l'intervalle $\icc{M}{\sqrt{2}}$, finalement, on a $M < \alpha < \sqrt{2}$ avec $\alpha \in \Q$ et donc, $M$ n'est pas un majorant.\+
   En conclusion, les majorants de $\E$ sont nécessairement dans $\ico{\sqrt{2}}{\infty}$.
\end{proof}

\subsubsection*{\subsubsecstyle{Cherchons l'ensemble des minorants ($MINO(\E$)) de $\E$:}}

L'intervalle équivalent à $\E$ étant symétrique, on peut trouver l'ensemble des minorants de manière analogue.

\subsubsection*{\subsubsecstyle{Max et Min de $\E$:}}

Le maximum (resp.\ minimum) de $\E$ est un majorant (resp.\ minorant) qui appartient à $\E$, or les majorants (resp.\ minorants) sont les réels $x$ tels que $x \geq \sqrt{2}$ (resp. $x \leq -\sqrt{2}$) et $\sqrt{2}$ (resp. $-\sqrt{2}$) n'appartiennent pas à $\E$ car ils sont irrationnels, donc $\E$ n'admet ni minimum, ni maximum.

\subsubsection*{\subsubsecstyle{Sup et Inf de $\E$:}}

sup$(\E)$= $\sqrt{2}$ car min$(MAJO(\E)) = \sqrt{2}$.

inf$(\E)$ = $-\sqrt{2}$ car max$(MINO(\E)) = -\sqrt{2}$.
\pagebreak

\subsection[Exercice 12:]{\subsecstyle{\textbf{Exercice 12:}}}
Soit \(A \subseteq \R\) une partie non vide et minorée, soit \(\alpha\) sa borne inférieure.\+
Soit \(B := A \cap \ioo{-\infty}{\alpha + 1}\).

\subsubsection*{\subsubsecstyle{(1) Justifier que \(\alpha\) existe:}}
\(A\) est une partie non vide et minorée de \(\R\) donc par la propriété de la borne supérieure, \(\alpha\) existe bel et bien.

\subsubsection*{\subsubsecstyle{(2) Justifier que B est non vide et minorée:}}
On sait que A est non vide, et que \(\alpha\) est le plus grand minorant de A. Supposons que B soit vide, cela signifie donc que \(A \cap \icc{-\infty}{\alpha + 1} = \emptyset\), or A est non vide donc on en déduis:
\[
   \lambda \in A \text{ ET } \lambda \notin \icc{-\infty}{\alpha + 1} \implies \lambda \geq \alpha + 1
\]
Et donc \(\alpha + 1\) serait un minorant de A plus grand que \(\inf(A)\) ce qui est absurde. Donc B est nécessairement non vide.\<

B est minorée car sachant que les éléments de B appartiennent à A on a:
\[
   \lambda \in B \implies \lambda \in A \implies \lambda \leq \alpha   
\]

\subsubsection*{\subsubsecstyle{(3) Justifier que B est non vide et minorée:}}
Par la question (2) on déduis que B admet une borne inférieure par la propriété de \(\R\). Et sachant que tout les éléments de B sont des éléments de A, par unicité de la borne inférieure, on conclus que \(\inf(B) = \alpha\).
\pagebreak

\subsection[Exercice 13:]{\subsecstyle{\textbf{Exercice 13:}}}
\begin{proof}[\unskip\nopunct]
   Soit \(A \subseteq \R\) une partie non vide et minorée telle que \(\sup(A) \notin A\).

   Montrons que pour tout \(\epsilon\) positif, \(I := \ioo{\sup(A) - \epsilon}{\sup(A)}\) contient une infinité d'éléments de A.\<

   Montrons tout d'abord la propriété intermédiaire \(\P\):
   \[
      \sup(A) \notin A \implies A \text{ n'est pas fini.}
   \]
   Prouvons \(\P\) par l'absurde et supposons que A soit un ensemble fini tel que:
   \[
      A := \Bigl\{ a_1, a_2 \ldots, a_{n-1}, a_{n} \; ; \; n \in \N \Bigl\} \text{ ET } \sup(A) \notin A
   \]
   Alors, par la relation d'ordre total de \(\R\), A admet un maximum, et on sait que si le maximum existe, il est la borne supérieure de l'ensemble, ce qui est absurde car par hypothèse, on a \(\sup(A) \notin A\), donc pas l'absurde, A est nécessairement infini (*).\<

   Montrons maintenant que I contient nécessairement une infinité d'éléments de A, on sait que A est minorée donc deux cas se présentent:
   \[
      \begin{cases}
         \epsilon < \inf(A) \implies A \subseteq I\\
         \epsilon \geq \inf(A) \implies I \subseteq A
      \end{cases}
   \]
   Dans le premier cas, on voit que A est inclus dans I, or A est infini (*), donc I contient une infinité d'éléments de A.

   Dans le second, cas on peut remarque que \(I\) est une partie de \(A\), or, en tant que partie de A elle ne contient toujours pas \(\sup(A)\), et donc par (*), on a que \(I\) contient une infinité d'élements, or \(I\) est une partie de A, donc tout ses éléments sont dans A.
\end{proof}
\pagebreak



